\documentclass[11pt]{article}
\usepackage[T1]{fontenc}
\usepackage{mathpazo}
\usepackage[margin=0.9in]{geometry}
\usepackage{amsmath,amssymb}
\usepackage{enumitem}
\usepackage[dvipsnames]{xcolor}
\usepackage{titlesec}
\usepackage{hyperref}
\hypersetup{colorlinks,urlcolor=blue,linkcolor=black}

\definecolor{oucrimson}{RGB}{132,22,23}
\titleformat{\section}{\large\bfseries\color{oucrimson}}{\thesection.}{0.5em}{}
\titlespacing*{\section}{0pt}{1.0ex}{0.4ex}
\setlist{nosep,leftmargin=1.3em}
\setlength{\parindent}{0pt}
\setlength{\parskip}{0.4ex}
\pagestyle{empty}

\newcommand{\review}[1]{\par\smallskip\textit{Review:} #1}

\begin{document}

{\LARGE\bfseries Midterm Exam Study Guide}\\[0.3ex]
{\large Econ 6343: Econometrics III, Fall 2026}

\smallskip
\fcolorbox{oucrimson}{oucrimson!6}{\parbox{\dimexpr\linewidth-2\fboxsep-2\fboxrule}{\centering\textbf{\color{oucrimson}Generated by Claude Opus 5.5}}}

\medskip
\textbf{When:} Tuesday, November 10, 2026, in class (75 minutes).\quad
\textbf{Format:} 7 questions, 70 points. Closed AI, notes, book, and device; a basic scientific calculator is allowed.\quad
\textbf{Coverage:} everything through Class 22 (Learning models).

Most questions ask you to \emph{explain}, \emph{critique a claim}, or decide whether a statement is \emph{true or false} and say why. Expect one short derivation and one question that asks you to read a short Julia function. Good answers are a few clear sentences, not long essays.

\section{Stated preferences and stated probabilities}
\begin{itemize}
\item Why researchers collect stated-preference data, and how they differ from revealed-preference data
\item What a stated probability measures, and what it adds beyond a discrete choice or a ranking
\item How logit-type models are estimated when the outcome is a probability rather than a 0/1 choice, and the practical issues that come up
\end{itemize}
\review{Lecture 11; Stated Preference videos; \emph{Train} 7.2--7.3; mobility (PCM) in-class activity.}

\section{Identification in dynamic discrete choice models}
\begin{itemize}
\item What it means for a parameter of a dynamic model to be identified, and which parameters are hardest to pin down
\item The role of functional-form assumptions, and what kinds of variation in the data help
\end{itemize}
\review{Lecture 7; Dynamic Discrete Choice videos (the whole series); Rust (1987).}

\section{Experiments, structural models, and counterfactuals}
\begin{itemize}
\item Internal vs.\ external validity; what experiments can and cannot tell us
\item How experiments and structural models can complement each other (estimation, validation, prediction)
\item What must be true for a counterfactual simulation to be credible
\end{itemize}
\review{Lectures 2, 3, and 10; Model Fit, Validation, and Counterfactuals videos; counterfactuals in-class activity.}

\section{Breadth across the course (short true/false items)}
\begin{itemize}
\item Interpreting the key parameters of nested and mixed logit models, and the restrictions theory places on them
\item Measuring goodness of fit in discrete choice models; consumer surplus and the logsum
\item Properties of simulation-based estimators (SML, SMM, indirect inference) and how they depend on the number of draws
\item What identifies preference heterogeneity; finite mixtures and the EM algorithm
\item Solving finite- vs.\ infinite-horizon dynamic models
\end{itemize}
\review{Lectures 4--7 and 9; the corresponding videos; \emph{Train} Ch.\ 3, 4, 6, and 10.}

\section{Reading and testing code}
\begin{itemize}
\item Read a short Julia function like those in the problem sets and in-class demos, and explain what each step computes
\item Connect the math of a model (probabilities, likelihoods, normalizations) to the lines of code that implement it
\item Spot mistakes that do \emph{not} produce an error message, and design unit tests that would catch them
\end{itemize}
\review{PS2--PS4 and PS7 (including the demo files); Clean Code lecture.}

\section{Conditional choice probabilities}
\begin{itemize}
\item The Hotz--Miller inversion: what it says and why it matters computationally
\item Terminal, renewal, and finite-dependence actions
\item What CCP methods do and do not simplify (estimation vs.\ counterfactuals)
\end{itemize}
\review{Lecture 8; CCP videos; Hotz and Miller (1993); Arcidiacono and Miller (2011); PS6 and its demo.}

\section{Unobserved heterogeneity: factor and learning models}
\begin{itemize}
\item Proxy variables, measurement error, and factor analysis (including normalizations and identification)
\item Learning models and Bayesian updating, and how beliefs evolve over time
\item How factor models and learning models differ, and how they relate to each other
\end{itemize}
\review{Lectures 12--13; Measurement Error and Learning Models videos; Heckman, Stixrud and Urzua (2006); Miller (1984); PS8.}

\section{Facts worth knowing cold}
\begin{itemize}
\item Standard Type 1 Extreme Value: mean $\gamma\approx 0.5772$ (Euler's constant), variance $\pi^2/6$; the difference of two is logistic
\item Logit choice probabilities $P_{ij} = \exp(v_{ij})/\sum_k \exp(v_{ik})$; only utility differences matter, and the scale is normalized
\item With i.i.d.\ T1EV errors, $\mathbb{E}\max_j\{v_j+\epsilon_j\} = \ln\sum_k\exp(v_k) + \gamma$
\end{itemize}

\section{Not on the exam}
Material from Class 24 onward: referee reports, the final-project demonstration, advanced optimization, extreme value theory beyond Type 1 EV, relative correlation restrictions, and machine learning. Presentation and referee-report content is also excluded.

\end{document}
